
\DocumentMetadata{tagging = off} 

\documentclass{ltx-talk}
\usepackage[fontset = fandol]{ctex}

%\setmainfont{TeX Gyre Termes}
%\setmathfont{STIX Two Math}




	

\EditInstance{footer}{std}{
	color         = black,
	separator     = \hfill|\hfill,
	element-order = {framenumber},
	%element-order    = {date, author, title, institute, framenumber},
	right-hspace  = 3mm,
	left-hspace   = 3mm,
}


%\usepackage{amsmath} 
%\newtheorem{definition}{Definition}
\definecolor{myblue}{rgb}{0.0, 0.1, 0.8}


% Preamble
\usepackage{amsthm}
% 在导言区添加
\newtheorem*{theorem}{Theorem}
\newtheorem*{corollary}{Corollary}
\newtheorem*{lemma}{Lemma}
\newtheorem*{definition}{Definition}
\newtheorem*{example}{Example}
\newtheorem*{remark}{Remark}   % 不编号的 remark
\newtheorem*{proposition}{Proposition}


\usepackage[ruled,vlined]{algorithm2e}
\SetNlSty{textbf}{\color{red}}{}







% 3. Define a simple command for Chinese text

%\newfontfamily\zhfont{FandolSong-Regular.otf}[
% Script=CJK,
% Renderer=HarfBuzz % Use the faster HarfBuzz engine in LuaLaTeX
%]
%\newcommand{\zh}[1]{{\zhfont #1}}


%\setmainfont{Times New Roman}

%\setmathfont{STIX Two Math}


%\usepackage[fontset=fandol]{ctex}

\pagecolor{green!5}



\title{the first presentation of Hazhy}
\subtitle{the first presentation of Hazhy}
\author{Hazhy}
\institute{None}
\date{2026-1-13}



\begin{document}
	\maketitle
	
	\begin{frame} 
		\frametitle{Sparsity and Manifolds}
		% English and Math stay in the default New Computer Modern Sans/Math
		The sparsity is defined by the $\ell_1$-norm:
		\[ \min \|x\|_{\ell_1} \quad \text{subject to} \quad Ax = b \]
		
		\begin{itemize}
			\item \textbf{Riemannian Manifold:} 
			\item \textbf{Optimization:} 
			\item  黑暗中海燕1
		\end{itemize}
		
		% Test the \ell symbol:
		The cursive $\ell$ is preserved from the default NewCMSansMath.
	\end{frame}

	\begin{frame}
	
		\begin{proposition}
			
		\end{proposition}
		\begin{definition}
			
		\end{definition}
		\begin{lemma}
			
		\end{lemma}
	\end{frame}
	
	
	\begin{frame*}{Sample Code}
		
		\begin{verbatim}
			print("Hello, sparse recovery!")
		\end{verbatim}
		
		
	\end{frame*}
	
	
	
	\begin{frame}
		\frametitle{中文测试}
		将军不下马，出入平安xasxxasxax
		你好i 撒CDC四川省成都成都vvvvvvvvvvvvvvvvvvvvvvvcvcvcvc
	\end{frame}
	
	\begin{frame}{Algorithm: Basis Pursuit}
		\begin{algorithm}[H]
			\caption{L1-Minimization}
			\KwIn{Matrix $A$, vector $b$}
			\KwOut{Sparse vector $x$}
			Solve $\min \|x\|_1$ subject to $Ax = b$\;
		\end{algorithm}
	\end{frame}
	
\end{document}
